% !TEX root = ../Main.tex
\chapter{仰泳}

\section{出发}

\state{出发前的准备}{
    \begin{itemize}
        \item 在出发信号发出前，运动员应\textbf{在水中面对出发端}，\textbf{两手抓住出发握手器}；
        \item \textbf{禁止}两脚蹬在水槽里、水槽上，或脚趾\textbf{勾在水槽沿上}。
    \end{itemize}
}

\section{游进姿势}

\state{仰卧要求}{
    \begin{itemize}
        \item 出发和每次转身后，运动员应\textbf{蹬离池壁}；
        \item 除做\textbf{转身动作}外，运动员在整个游进过程中\textbf{应始终呈仰卧姿势}；
        \item 允许身体\textbf{转动}，但必须保持\textbf{与水平面小于 90$^\circ$ 的仰卧姿势}；
        \item \textbf{头部位置不受此限}（即头部可以不严格保持仰卧）。
    \end{itemize}
}

\state{15 米规则（仰泳版本）}{
    \begin{itemize}
        \item 整个游进中，\textbf{身体的某一部分必须露出水面}；
        \item 但在\textbf{出发、转身及抵达终点}时，允许身体\textbf{完全没入水中}；
        \item 出发和每次转身后，在\textbf{15 米前（含 15 米）头部必须露出水面}。
    \end{itemize}
}

\section{转身与终点}

\state{转身动作}{
    \begin{itemize}
        \item 转身过程中，\textbf{身体的某一部分必须触壁}；
        \item 允许\textbf{肩的转动超过垂直面}；
        \item 之后可做\textbf{一次单臂划水}或\textbf{双臂同时划水}动作，以此作为连贯转身动作的开始；
        \item 运动员必须\textbf{呈仰卧姿势蹬离池壁}。
    \end{itemize}
}

\state{终点触壁}{
    运动员在到达终点时，\textbf{必须以仰卧姿势触壁}。
}

\rmkb{
    \textbf{易考点}：
    \begin{itemize}
        \item \textbf{仰泳在水中出发}，两手抓\textbf{握手器}、两脚\textbf{不得蹬水槽或勾沿}；
        \item 整个游程\textbf{始终仰卧}（仅转身动作除外），\textbf{头部位置不限}；
        \item 转身后必须\textbf{仰卧蹬壁}；
        \item \textbf{终点必须仰卧触壁}；
        \item 15 米规则同自由泳。
    \end{itemize}
}
